\section{Sheaves on $[\Spec A / G]$} \label{section: sheaves on [spec A /G]}

In the previous section, we introduced a stacky curve $\mc{X}$ and studied its fundamental group. In what follows, we want to study different moduli spaces of bundles on $\mc{X}$ and understand how they are related to the corresponding moduli spaces on $X$. As observed in \eqref{X|U}, the local structure of $\mc{X}$ around the stacky point is $[\Ab^1/\mu_e]$. Hence, a first step for understanding bundles on $\mc{X}$ is to understand coherent sheaves on $[\Ab^1 / \mu_e]$. In this section, we study more in general coherent sheaves on quotient stacks of the form $[\Spec A / G]$, where $G$ is a discrete group acting on $\Spec A$.

In the rest of the section, we denote $\Spec A$ by $X$ and we work over $\Spec k$, where $k$ is an algebraically closed field of characteristic $0$.

\subsection{Coherent sheaves}\label{coh([X/G])}

The quotient map $ X \to [ X / G]$ is an atlas. Therefore, the datum of coherent sheaf on $[ X/ G]$ is equivalent to the datum of a coherent sheaf $\Fc$ on $X $ with descent data. The descent data is composed by an isomorphism
\[
\alpha \colon pr_1^* \Fc \to pr_2^* \Fc
\]
of coherent sheaves on $X \times_{BG} X$, and by natural isomorphisms on $X\times_{BG}X\times_{BG} X$, giving the cocycle condition.

\begin{rmk}
There is an isomorphism $X \times_{BG} X \cong G \times X$. Composing with this isomorphism, the projection $pr_2 \colon  X \times_{BG} X \to X$ is the natural projection $ G \times X \to X$ and the projection $pr_1 \colon X \times_{BG} X \to X $ is the group action $\sigma \colon G \times X \to X$.
\end{rmk}

In the case $X = \Spec A$, the coherent sheaf $\mc{F}$ corresponds to a finitely presented $A$-module $M$. The isomorphism of $\O_{G \times X}$-modules
\[
\alpha \colon  pr_2^* \Fc \to \sigma^* \Fc
\]
corresponds to an isomorphism of $(k[G] \otimes A)$-modules
\[
\alpha^{\#} \colon  k[G] \otimes M \to k[G] \otimes M,
\]
where on the left hand side we have the natural structure of $k[G] \otimes A$-module, while on the right hand side $k[G]$ acts also on $M$ via the action on $A$.
This gives a morphism of $A$-modules
\[
\rho \colon M \to k[G] \otimes M.
\]
The morphism $\rho$ gives $M$ the structure od a $k[G]$-comodule. Indeed, the commutativity of
% \[\begin{tikzcd}
% 	{k[G]} & {k[G] \otimes V} \\
% 	{k[G] \otimes V} & {k[G] \otimes k[G] \otimes V}
% 	\arrow["\rho"', from=1-1, to=1-2]
% 	\arrow["\rho", from=1-1, to=2-1]
% 	\arrow["{\id_k[G] \otimes \rho}", from=1-2, to=2-2]
% 	\arrow["{\Delta \otimes \id_V}"', from=2-1, to=2-2]
% \end{tikzcd}\]
\[\begin{tikzcd}
	{k[G]} && {k[G] \otimes M} \\
	{k[G] \otimes M} && {k[G] \otimes k[G] \otimes M}
	\arrow["\rho", between={0.1}{0.9}, from=1-1, to=1-3]
	\arrow["\rho", from=1-1, to=2-1]
	\arrow["{\id_{k[G]} \otimes \rho}", from=1-3, to=2-3]
	\arrow["{\Delta \otimes \id_M}"', from=2-1, to=2-3]
\end{tikzcd}\]
follows from the cocycle condition expressed by
\[\begin{tikzcd}
	{X \times_{[X/G]} X \times_{[X/G]} X } & {G \times G \times X } && {G \times X} & X.
	\arrow[between={0.3}{0.7}, equals, from=1-1, to=1-2]
	\arrow["{\id_G \times \sigma}", from=1-2, to=1-4]
	\arrow["{m \times \id_X}", shift left, curve={height=-12pt}, from=1-2, to=1-4]
	\arrow["{pr_{13}}", shift right, curve={height=12pt}, from=1-2, to=1-4]
	\arrow["{pr_2}"', shift right=2, from=1-4, to=1-5]
	\arrow["\rho", shift left=2, from=1-4, to=1-5]
\end{tikzcd}\]
While the commutativity of
\[\begin{tikzcd}
	M & {k[G] \otimes M} \\
	& {M \cong k \otimes M}
	\arrow["\rho", from=1-1, to=1-2]
	\arrow["{\id_M}"', shift right, curve={height=12pt}, from=1-1, to=2-2]
	\arrow["{\varepsilon \otimes \id_M}", from=1-2, to=2-2]
\end{tikzcd}\]
follows from the fact that $\alpha^{\#}$ is an isomorphism of $k[G]$-modules.

In conclusion, the datum of a coherent sheaf on $[\Spec A / G]$ is equivalent to the datum of an $A$-module $M$ with the structure of $k[G]$-comodule. But to give $M$ the structure of $k[G]$-comodule is equivalent to giving it the structure of a $G$-representation. Therefore, coherent sheaves on $[\Spec A / G]$ are $A$-modules with a $G$-action. In the case $A = k$, we get an equivalence
\[
\left< \text{Coherent sheaves on } BG \right> \cong \left< \text{Finite dimensional } G\text{-representations} \right>.
\]

\subsection{Vector bundles}

As just mentioned, a coherent sheaf on $BG$ gives a $G$-representation. Let us consider more in detail the case of vector bundles.
A vector bundle $\Ec$ on $BG$ corresponds to a morphism $BG \to BGL_n$.

\begin{rmk}
A morphism $f \colon BG \to BGL_n$ yields a homomorphism of group schemes $G \to GL_n$.
\end{rmk}
\begin{proof}
    For any test scheme $T$, consider the trivial $G$-torsor $G \times T \to T$ as an object of $BG(T)$. Then
    \[
    \Aut_{BG(T)}(G \times T \to T) = G(T).
    \]
    Notice that $f$ must map trivial torsors to trivial torsors.
	Indeed, for $T = \Spec k$ we have that $f(G \to \Spec k)$ is a $GL_n$-torsor over $\Spec k$. Since $k$ is algebraically closed, $\Spec k$ admits only trivial torsors, therefore $f(G \to \Spec k)$ is the trivial torsor $GL_n \to \Spec k$. For other schemes $T \in \Sch/k$ the statement follows by compatibility with pullback.\\
	% Indeed, this is trivial for $T = \Spec k$ and follows by compatibility under pullback for other schemes $T \in \Sch/k$.\\
    % (In this part of the argument it is essential that $k$ is algebraically closed, so that $\Spec k$ admits only the trivial torsor)\\
    Passing to the automorphisms groupoids, we get that $f$ induces a morphism of groupoids
    \[
    G(T) = \Aut_{BG(T)}(G \times T \to T) \to \Aut_{BGL_n}(GL_n \times T \to T ) = GL_n(T),
    \]
    which is functorial in $T$. This is exactly a morphism of group schemes $G \to GL_n$, i.e. an $n$-dimensional representation of $G$.
\end{proof}

% \begin{rmk}
% Notice that the representation obtained by the above construction comes from the descent datum associated to the vector bundle $\Ec$ on $BG$. For this reason, we get exactly the same representation as in \zcref{coh([X/G])}.
% \end{rmk}

\begin{rmk}
The $G$-representation $G \to GL_n$ found above is the same as the representation obtained by the descent data of the vector bundle $\mc{E}$ corresponding to the morphism $BG \to BGL_n$, as mentioned in \zcref{coh([X/G])}.
\end{rmk}


\subsection{The case of a stacky curve} \label{subsection: the case of a stacky curve}
% Let us now focus on the case we are interested in. Let $X$ be a complex curve and $p \in X$ be a complex point. Associated to $p$ is a pair $(\O(p), s_p)$, where $O(p)$ is the line bundle associated to the effective Cartier divisor $p$ and $s_p$ is the \textit{tautological section} of $\O(p)$, vanishing at $p$.

% The \textit{stacky curve} $\Xc$ is defined as the pullback
% % \[\begin{tikzcd}
% % 	\Xc & {[\Ab^1/\Gb_m]} \\
% % 	X & {[\Ab^1/\Gb_m]}
% % 	\arrow[from=1-1, to=1-2]
% % 	\arrow[from=1-1, to=2-1]
% % 	\arrow["\lrcorner"{anchor=center, pos=0.125}, draw=none, from=1-1, to=2-2]
% % 	\arrow["{\wedge n}", from=1-2, to=2-2]
% % 	\arrow["{(\O(p), s_p)}"', from=2-1, to=2-2]
% % \end{tikzcd}\]
% \[\begin{tikzcd}
% 	\Xc && {[\Ab^1/\Gb_m]} \\
% 	X & {} & {[\Ab^1/\Gb_m].}
% 	\arrow[from=1-1, to=1-3]
% 	\arrow[from=1-1, to=2-1]
% 	\arrow["\lrcorner"{anchor=center, pos=0.125}, draw=none, from=1-1, to=2-2]
% 	\arrow["{\wedge e}", from=1-3, to=2-3]
% 	\arrow["{(\O(p), s_p)}"', from=2-1, to=2-3]
% \end{tikzcd}\]
% The stack $\Xc$ comes equipped with a pair $(\O(\frac{1}{e}p), \frac{1}{e}s_p)$, where $\O(\frac{1}{e}p)$ is a line bundle and $\frac{1}{e}s_p$ is a section. This pair corresponds to the upper horizontal morphism, and is characterized by the properties
% \[
% \pi^*\O(p) = \O(\tfrac{1}{e}p)^e, \quad \pi^* s_p = (\tfrac{1}{e}s_p)^e.
% \]
% The sheaf $\O(\frac{1}{e}p)$ is called the \textit{tautological sheaf} of the stacky curve $\Xc$. It corresponds to the composition
% \[
% \Xc \to [\Ab^1 / \Gb_m] \to B\Gb_m.
% \]

We now apply the theory to the example presented in \zcref{stacky curve}.
As already mentioned, $\mc{X}$ is equipped with a line bundle $\O(\tfrac{1}{e}p)$.

We are interested in studying the behavior of $\O(\frac{1}{e}p)$ at the stacky point of $\Xc$. Namely, considering the fiber product
\[\begin{tikzcd}
	{B\mu_e} & \Xc \\
	{*}      & X,
	\arrow["i", from=1-1, to=1-2]
	\arrow[from=1-1, to=2-1]
	\arrow["\lrcorner"{anchor=center, pos=0.125}, draw=none, from=1-1, to=2-2]
	\arrow["\pi", from=1-2, to=2-2]
	\arrow["p", from=2-1, to=2-2]
\end{tikzcd}\]
we are interested in the pullback $i^* \O(\frac{1}{e}p)$, which corresponds to the composition
\[
B\mu_n \xrightarrow{i} \Xc \to [\Ab^1 / \Gb_m] \to B\Gb_m.
\]

\begin{thm}\cite[Theorem 1.1.36]{Taams}
In the following diagram, the outer square is cartesian
\[\begin{tikzcd}
	{B\mu_e} & \Xc & {[\Ab^1/\Gb_m]} & {B\Gb_m} \\
	{*} & X & {[\Ab^1/\Gb_m]} & {B\Gb_m.}
	\arrow[from=1-1, to=1-2]
	\arrow[from=1-1, to=2-1]
	\arrow["\lrcorner"{anchor=center, pos=0.125}, draw=none, from=1-1, to=2-2]
	\arrow[from=1-2, to=1-3]
	\arrow["\pi", from=1-2, to=2-2]
	\arrow["\lrcorner"{anchor=center, pos=0.125}, draw=none, from=1-2, to=2-3]
	\arrow[from=1-3, to=1-4]
	\arrow["{\wedge e}", from=1-3, to=2-3]
	\arrow["{\wedge e}", from=1-4, to=2-4]
	\arrow["p"', from=2-1, to=2-2]
	\arrow[from=2-2, to=2-3]
	\arrow[from=2-3, to=2-4]
\end{tikzcd}\]
In particular, it is the pullback square
\[\begin{tikzcd}
	{B\mu_e} & {B\Gb_m} \\
	{*} & {B\Gb_m}
	\arrow[from=1-1, to=1-2]
	\arrow[from=1-1, to=2-1]
	\arrow["\lrcorner"{anchor=center, pos=0.125}, draw=none, from=1-1, to=2-2]
	\arrow["{\wedge e}", from=1-2, to=2-2]
	\arrow[from=2-1, to=2-2]
\end{tikzcd}\]
corresponding to the exact sequence of groups
\[
1 \to \mu_e \to \Gb_m \xrightarrow{\wedge e} \Gb_m.
\]
\end{thm}
\begin{proof}
    Consider the left side of the diagram
    \begin{equation}\label{cd}
    \begin{tikzcd}
	{B\mu_e} &&&& {[\Ab^1/\Gb_m]} \\
	{*} &&&& {[\Ab^1/\Gb_m].}
	\arrow["{{(\pi^*\O(p)\vert_{B\mu_e}, \; \pi^*s_p\vert_{B\mu_e})}}", from=1-1, to=1-5]
	\arrow["\psi", from=1-1, to=2-1]
	\arrow["{{{\wedge e}}}", from=1-5, to=2-5]
	\arrow[""{name=0, anchor=center, inner sep=0}, "{{(\O(p)\vert_p, \; s_p\vert_p)}}"', from=2-1, to=2-5]
	\arrow["\lrcorner"{anchor=center, pos=0.125}, draw=none, from=1-1, to=0]
    \end{tikzcd}
    \end{equation}
    Since the section $s_p$ by definition vanishes at $p$, we have $s_p \vert_p = 0$. Moreover, by construction
    \[
    (\pi^* s_p \vert_{B\mu_e})^e = \psi^* s_p \vert_p = \psi^* 0 = 0,
    \]
    therefore $\pi^* s_p \vert_{B\mu_e}$ = 0. This shows that, in the diagram \eqref{cd}, the horizontal arrows factor through the inclusion $B\Gb_m \to [\Ab^1 / \Gb_m]$ given by the choice of the zero section. So we get
    \begin{equation}\label{cd2}
    \begin{tikzcd}
	{B\mu_e} && {B\Gb_m} & {[\Ab^1/\Gb_m]} \\
	{*} && {B\Gb_m} & {[\Ab^1/\Gb_m],}
	\arrow["{\O(\frac{1}{e}p) \vert_B\mu_e}", from=1-1, to=1-3]
	\arrow[from=1-1, to=2-1]
	\arrow[from=1-3, to=1-4]
	\arrow[dashed, from=1-3, to=2-3]
	\arrow["{{{\wedge e}}}", from=1-4, to=2-4]
	\arrow["{\O(p)\vert_p}", from=2-1, to=2-3]
	\arrow[from=2-3, to=2-4]
    \end{tikzcd}
    \end{equation}
    where the vertical arrow can be filled by the map induced by the $e$-th power morphism $\Gb_m \to \Gb_m$. Then, the right square in \eqref{cd2} becomes the cartesian square
    \[\begin{tikzcd}
	{B\Gb_m} & {[\Ab^1/\Gb_m]} \\
	{B\Gb_m} & {[\Ab^1/\Gb_m].}
	\arrow[from=1-1, to=1-2]
	\arrow["{\wedge e}", from=1-1, to=2-1]
	\arrow["\lrcorner"{anchor=center, pos=0.125}, draw=none, from=1-1, to=2-2]
	\arrow["{{{\wedge e}}}", from=1-2, to=2-2]
	\arrow[from=2-1, to=2-2]
    \end{tikzcd}\]
    This shows that in \eqref{cd2} both the outer and the rightmost squares are cartesian. Therefore, the leftmost square is also cartesian. This means that $\O(\frac{1}{e}p)\vert_{B\mu_e}$ is the line bundle obtained by the pullback square
    \[\begin{tikzcd}
	{B\mu_e} & {B\Gb_m} \\
	{*} & {B\Gb_m,}
	\arrow[from=1-1, to=1-2]
	\arrow[from=1-1, to=2-1]
	\arrow["\lrcorner"{anchor=center, pos=0.125}, draw=none, from=1-1, to=2-2]
	\arrow["{\wedge e}", from=1-2, to=2-2]
	\arrow[from=2-1, to=2-2]
    \end{tikzcd}\]
    corresponding to the inclusion $\mu_e \hookrightarrow \Gb_m$.
\end{proof}

\begin{rmk}
The one above is a modified version of the proof given in \cite[]{Taams}, which we think might provide more clarity.
\end{rmk}

This means that the $\mu_e$-representation obtained as $\O(\frac{1}{e}p)\vert_{B\mu_e}$ is just the tautological character, corresponding to the inclusion $\mu_e \hookrightarrow \Gb_m$.